\section{Related Work}
\label{app:related-work}

Approaches to efficient reasoning broadly include test-time compute allocation, RL-based optimization, and distillation. 

\paragraph{Test-time compute allocation.}
At inference time, reasoning computation can be adjusted through the length of reasoning traces, the number of sampled solutions, and verification or search. Chain-of-thought prompting and self-consistency illustrate the benefits of explicit reasoning and exploring multiple solution paths~\citep{wei2022chain,wang2022self}, while process supervision supports verifier-guided selection~\citep{lightman2024lets}. However, longer reasoning can waste computation on simple problems~\citep{chen2024not}. \citet{snell2024scaling} show that effective compute allocation depends on prompt difficulty, motivating adaptive inference strategies. \citet{wan2026shadow} further formulate cross-query budget allocation using a global shadow price to balance the marginal utility of reasoning computation. Budget forcing provides another way to control the amount of reasoning at test time~\citep{muennighoff2025s1}. These methods adjust how much computation a model spends when answering a query.

\paragraph{RL-based reasoning efficiency.}
RL can train policies to use computation more efficiently. L1 optimizes accuracy together with adherence to requested reasoning-length constraints~\citep{aggarwal2025l1}, and lazy length penalties incorporate response-length reduction into reasoning RL~\citep{yuan2026shorten}. DAST uses difficulty-dependent budgets, reward shaping, and preference optimization to discourage excessive reasoning on easier problems while retaining sufficient computation for harder ones~\citep{shen2025dast}. AdapThink adapts length penalties to query difficulty~\citep{xu2026adapthink}. Training efficiency can also be improved through sample selection: DAPO filters rollout groups with uninformative rewards~\citep{yu2025dapo}, while online difficulty filtering focuses learning on tasks of intermediate difficulty~\citep{bae2026online}. Under group-relative objectives, however, all-correct groups have no reward-based relative advantage~\citep{shao2024deepseekmath}. This leaves little direct signal for preserving their existing concise behavior as the policy continues learning from other prompts.

\paragraph{Distillation and LSD.}
Distillation provides token-level supervision for learning concise reasoning. Generalized Knowledge Distillation trains students on their own generated sequences using teacher feedback, supports different divergence objectives, and can be combined with RL fine-tuning~\citep{agarwal2024onpolicy}. CRISP uses a periodically refreshed student copy conditioned on a conciseness instruction as its teacher and optimizes reverse KL on all student rollouts~\citep{sang2026opsdc}. Contrastive On-Policy Distillation compares teacher probabilities under light- and heavy-reasoning instructions to construct token-level advantages~\citep{ruan2026contrastive}.

Unlike these approaches, LSD combines on-policy distillation with online difficulty routing. Its focus is preserving concise behavior on already-solved queries during ongoing RL training, which we evaluate through LST on fixed easy-query sets.
